\chapter*{Resumen}

    La vista es uno de los sentidos que nos ayudan a interpretar el entorno y a interactuar de manera efectiva con él. Sin embargo, personas que carecen de manera total o parcial de este sentido se enfrentan a desafíos diarios que limitan su independencia y acceso a diversas áreas cotidianas. La navegación para personas cuya visión es reducida o presentan discapacidad visual es esencial para su autonomía y bienestar. Si bien existen métodos tradicionales de asistencia, estos suelen tener un enfoque  que no se ajuste a las necesidades de viajar a diversos puntos. Este trabajo terminal tiene como objetivo utilizar la inteligencia artificial para proporcionar asistencia auditiva en tiempo real a personas con visión reducida, mediante el uso de lentes adaptados, así como un sistema de navegación. El sistema permite la identificación del entorno para guiar al usuario de manera adecuada, promoviendo una mayor independencia y seguridad para los usuarios.\\
        
    Para lograrlo, se implementa un sistema inteligente capaz de percibir y analizar el entorno, identificando elementos relevantes como obstáculos, señales y personas. Esta información se integra con datos provenientes de distintos sensores para construir una representación dinámica del espacio circundante. A partir de ello, el sistema es capaz de estimar la ubicación del usuario y generar rutas de navegación que se ajustan continuamente ante cambios en el entorno. Asimismo, se incorpora un módulo de análisis contextual que permite emitir alertas oportunas, junto con mecanismos de interacción natural que facilitan la comunicación con el usuario y hacen al sistema fácil de comprender. Finalmente, la información procesada se transforma en indicaciones claras y accesibles, conformando una solución integral de asistencia a la navegación.
    

\newpage

 \chapter*{Abstract}

    Sight is one of the senses that helps us interpret our surroundings and interact effectively with them. However, people who are totally or partially blind face daily challenges that limit their independence and access to various areas of everyday life, such as mobility. Navigation for people with low vision or visual impairments is essential for their autonomy and well-being. While traditional methods of assistance exist, they lack an accessible and efficient approach that meets the needs of traveling to various locations. This final project aims to use artificial intelligence to provide real-time auditory assistance to people with low vision, through the use of adapted lenses, an optimized navigation system, and location assistance in emergency situations. The system allows for the identification of the environment to guide the user appropriately, promoting greater independence and safety for users.

    To achieve this, an object detection model based on convolutional neural networks is implemented to identify obstacles, traffic signs, and people in the user’s surroundings in real time. This information is combined with depth sensor data through sensor fusion, generating a dynamic map of the environment using simultaneous localization and mapping techniques. Navigation decisions are made by a deep reinforcement learning agent (PPO) that guides the user step by step towards places recognized in the map, such as doors, and reacts immediately to unforeseen obstacles. A probabilistic inference module evaluates the context to issue relevant alerts, while a voice recognition system enables natural user interaction with the device. All processed information is translated into clear and timely audio instructions, forming a comprehensive assisted navigation system.


